\section{Kết luận và hướng phát triển}

\subsection{Kết luận}

Nghiên cứu này trình bày một khảo sát thực nghiệm có hệ thống về bốn chiến lược ngưỡng quyết định thích nghi - cố định, phân bin, tuyến tính, và tương tác - cho bài toán nhận diện khuôn mặt trong điều kiện ánh sáng biến đổi, hướng đến triển khai trên thiết bị biên. Thực nghiệm được tiến hành trên tập dữ liệu tự thu thập gồm 14 người với 1.144 ảnh ở ba điều kiện ánh sáng, sử dụng ArcFace (InsightFace \textit{buffalo\_sc}) làm backbone trích xuất đặc trưng.

Đối chiếu với ba câu hỏi nghiên cứu đặt ra ở Mục~1, chúng tôi rút ra các kết luận sau:

\begin{enumerate}
    \item \textbf{H1 - Hiệu quả của ngưỡng thích nghi: Xác nhận một phần.} Chiến lược ngưỡng phân bin (F2) giảm FRR ở điều kiện \texttt{dark} từ 9,8\% xuống 3,9\% (giảm 60\%, $\Delta\text{FRR} = -5{,}9$ pp), vượt mục tiêu đặt ra, trong khi FAR tăng ở mức chấp nhận được ($+1{,}7$ pp). Tuy nhiên, công thức tương tác (F4) không phát huy tác dụng ($\gamma^* = 0{,}0$) trên tập dữ liệu hiện tại, cho thấy mô hình phức tạp hơn không nhất thiết tốt hơn khi dữ liệu hiệu chỉnh có quy mô nhỏ.

    \item \textbf{H2 - Ổn định khi cập nhật gallery: Chưa kiểm chứng.} Do giới hạn về thời gian thực hiện, cơ chế cập nhật gallery trực tuyến và metric CP-BWT chưa được đánh giá trên dữ liệu thực tế. Đây là câu hỏi nghiên cứu mở cần được giải quyết trong giai đoạn tiếp theo.

    \item \textbf{H3 - Khả năng triển khai biên: Đạt về latency, chưa đạt về RAM.} Latency trung bình 108,6~ms (P95 = 118,6~ms) đáp ứng yêu cầu thời gian thực dưới 200~ms với biên an toàn đáng kể. Tuy nhiên, RAM peak 3,2~GB vượt xa mục tiêu 512~MB, đòi hỏi giải pháp tối ưu bộ nhớ (lượng tử hóa mô hình hoặc thay thế backbone) trước khi triển khai trên thiết bị tài nguyên hạn chế.
\end{enumerate}

\subsection{Đóng góp của nghiên cứu}

Tổng hợp lại, nghiên cứu này có ba đóng góp chính vào lĩnh vực nhận diện khuôn mặt thích nghi:

\begin{enumerate}
    \item \textbf{Khung so sánh thực nghiệm}: Cung cấp nghiên cứu so sánh có hệ thống đầu tiên (theo hiểu biết của chúng tôi) giữa bốn dạng hàm ngưỡng quyết định với embedding ArcFace thực tế trong bối cảnh triển khai chấm công tại Việt Nam. Kết quả cho thấy chiến lược phân bin đơn giản đạt hiệu quả cao nhất trên tập dữ liệu quy mô nhỏ, cung cấp hướng dẫn thực tiễn cho các hệ thống triển khai tương tự.

    \item \textbf{Metric CP-BWT}: Đề xuất \textit{Condition-Partitioned Backward Transfer} như một công cụ đo lường tính ổn định của hệ thống nhận diện theo phân vùng điều kiện môi trường, mở rộng khái niệm BWT truyền thống trong học liên tục sang lĩnh vực triển khai thực tế.

    \item \textbf{\textit{Negative finding} có giá trị}: Xác định rằng thành phần tương tác giữa độ sáng và độ nhiễu không phát huy tác dụng trên tập dữ liệu quy mô nhỏ, định hướng cho nghiên cứu tiếp theo về điều kiện cần thiết (quy mô dữ liệu, tính đại diện) để các mô hình ngưỡng phức tạp có thể vượt trội so với phương pháp đơn giản.
\end{enumerate}

\subsection{Giới hạn của nghiên cứu}

Nghiên cứu này có bốn giới hạn chính cần được lưu ý khi diễn giải kết quả:

\begin{enumerate}
    \item \textbf{Quy mô và tính đại diện của dữ liệu}: Tập dữ liệu chỉ gồm 14 người trong một môi trường duy nhất. Kết quả, đặc biệt là sự vượt trội của chiến lược Bin so với Linear và Interaction, có thể thay đổi khi quy mô tăng lên và phân bố dữ liệu phong phú hơn.

    \item \textbf{Dữ liệu \texttt{dark} tổng hợp}: Toàn bộ ảnh \texttt{medium} và \texttt{dark} được tạo bằng augmentation, chưa phản ánh đầy đủ sự phức tạp của điều kiện ánh sáng yếu tự nhiên. Điều này ảnh hưởng trực tiếp đến khả năng đánh giá công bằng công thức tương tác (F4) vốn được thiết kế cho tương tác thực tế giữa ánh sáng và nhiễu.

    \item \textbf{RAM chưa đạt mục tiêu}: Với 3,2~GB RAM peak, hệ thống chưa triển khai được trên các thiết bị biên phổ biến (Raspberry Pi 4B) mà không có tối ưu bổ sung.

    \item \textbf{H2 và CP-BWT chưa được đánh giá thực nghiệm}: Câu hỏi nghiên cứu H2 và metric CP-BWT mới dừng ở mức đề xuất lý thuyết, chưa có kết quả thực nghiệm để kiểm chứng.
\end{enumerate}

\subsection{Hướng phát triển}

Dựa trên các giới hạn đã xác định, chúng tôi đề xuất năm hướng phát triển theo thứ tự ưu tiên:

\begin{enumerate}
    \item \textbf{Thu thập dữ liệu \texttt{dark} thực tế}: Xây dựng tập dữ liệu với ít nhất 30--50 người, sử dụng camera thiết bị biên ghi nhận trực tiếp trong điều kiện ca đêm và hành lang thiếu sáng. Đây là điều kiện tiên quyết để kiểm chứng lại giả thuyết về hiệu ứng tương tác (F4) và đánh giá tính tổng quát của kết quả hiện tại.

    \item \textbf{Tối ưu bộ nhớ cho triển khai biên}: Lượng tử hóa mô hình InsightFace sang INT8 bằng ONNX Runtime hoặc TensorRT, kết hợp đánh giá tác động lên accuracy. Thay thế backbone bằng MobileFaceNet nếu yêu cầu RAM nghiêm ngặt ($<$ 512~MB).

    \item \textbf{Triển khai và đánh giá H2}: Cài đặt cơ chế cập nhật gallery trực tuyến, sử dụng metric CP-BWT để đo lường tính ổn định theo từng phân vùng điều kiện trước và sau khi cập nhật.

    \item \textbf{Kiểm chứng trên phần cứng thực tế}: Triển khai pipeline hoàn chỉnh trên Raspberry Pi 5 hoặc các SBC tương đương, đo latency và memory profile trong điều kiện vận hành liên tục (24h).

    \item \textbf{Mở rộng bài toán}: Tích hợp khả năng \textit{open-set recognition} (phát hiện người lạ không có trong gallery) và đánh giá hiệu năng của ngưỡng thích nghi trong ngữ cảnh này, nơi phân biệt giữa genuine pair có similarity thấp (do ảnh tối) và impostor pair trở nên khó khăn hơn.
\end{enumerate}